\documentclass[12pt]{article}
\usepackage[utf8]{inputenc}
\usepackage{amsmath, amssymb, amsthm}
\usepackage{geometry}
\usepackage{fancyhdr}

\geometry{letterpaper, margin=1in}
\setlength{\headheight}{15pt}

% Header setup
\pagestyle{fancy}
\fancyhf{}
\lhead{Math 292 Spring 2026}
\rhead{Jeffrey Lu}
\cfoot{\thepage}

\title{Homework Assignment 2 \\ \large Math 292, Spring 2026 \\ Rutgers University}
\author{Jeffrey Lu}
\date{2/13/2026}

\begin{document}

\maketitle

\noindent \textbf{Exercise 2.2:}

1. We can solve for the general solution by separation of variables
\begin{align*}
  x'(t)          & = tan(x)   \\
  cot(x)dx       & = dt       \\
  \int{cot(x)dx} & = \int{dt} \\
  ln|sin(x)|     & = t + C    \\
  sin(x)         & = Ce^t     \\
\end{align*}

2. We now solve for the constant C
\begin{align*}
  sin(x_0) & = Ce^{0}   \\
  C        & = sin(x_0)
\end{align*}

3. Plug back in for the particular solution
\begin{align*}
  sin(x) & = e^t sin(x_0)         \\
  x(t)   & = arcsin(e^t sin(x_0))
\end{align*}

4. With the given bounds of tan(x), we have that  $arcsin(e^tsin(x_0)) \neq{} \pm{\dfrac{\pi{}}{2}}$ and thus
\[
  -1 < e^t sin(x_0) < 1
\]

5. Since $e^t$ is always positive we can rewrite with absolute values
\begin{align*}
  e^t |sin(x_0)| & < 1            \\
  e^t            & < |csc(x_0)|   \\
  t              & < ln|csc(x_0)|
\end{align*}

6. Therefore, we can conclude the solution is defined when
\[
  t \in{(-\infty{},~ ln|csc(x_0)|)} ~|~ x_0 \neq{0}
\]

7. However, there is also the special case when $x_0=0$. In this case the solution no longer depends on t and thus
\[
  t\in{(-\infty, \infty)} ~|~ x_0 = 0
\]

\noindent \textbf{Exercise 2.3:}

\begin{enumerate}
  \item [(a)]

        1. First we plug in the initial condition and conclude that x is "increasing"
        \[
          v(0) = (1-0)^2 = 1
        \]

        2. Now we can calculate Barrow's Formula from x = 0 to x = 1
        \begin{align*}
          t(1) & = t(0) + \int^1_0{\frac{1}{v(z)} ~ dz}  \\
               & = \int^1_0{\frac{1}{\sqrt{1-z^4}} ~ dz} \\
        \end{align*}

        3. Now need to determine whether the integral converges or diverges. We can do this via the limit comparison test, the comparison integral $g(x)$ we choose is
        \begin{align*}
          g(x) & = \int^1_0{\frac{1}{\sqrt{1-z}} dz}                           \\
               & = \lim_{t\rightarrow{1^{-}}}\int^t_0{\frac{1}{\sqrt{1-z}} dz} \\
               & = \lim_{t\rightarrow{1^{-}}}[-2\sqrt{1-z}]^t_0                \\
               & = \lim_{t\rightarrow{1^{-}}} -2\sqrt{1-t} + 2                 \\
               & = 2
        \end{align*}

        4. Therefore, g(x) converges since its limit exists and is finite. Notice that the function becomes undefined when z approaches 1 from 0, that is why we are using the limit as it approaches 1 from the left. Now, we solve for the limit of the ratio $\frac{t(1)}{g(z)}$.
        \begin{align*}
          L & = \lim_{z\rightarrow{1^{-}}}\frac{\frac{1}{\sqrt{1-z^4}}}{\frac{1}{\sqrt{1-z}}}  \\
            & =  \lim_{z\rightarrow{1^{-}}}\frac{\sqrt{1-z}}{\sqrt{1-z}\sqrt{1+z}\sqrt{1+z^2}} \\
            & =                    \lim_{z\rightarrow{1^{-}}}\frac{1}{\sqrt{1+z}\sqrt{1+z^2}}  \\
            & = \frac{1}{2}
        \end{align*}

        5. L is both positive and finite, and thus we can conclude by the limit comparison test that t(1) converges. Therefore, the solution does not remain only within (-1,1) since it reaches the boundary 1.

        \vspace{20pt}

  \item [(b)]

        1. Taking the derivative of v(x), we get
        \begin{align*}
          v'(x) = -8x^3(1-x^4)
        \end{align*}

        2. This is a polynomial and when you plug in 1 or -1, you get 0. Thus, the function is continuous. Since v'(x) is continuous, we can conclude that the function is lipschitz and thus solution exist for all t and remain within (-1, 1)
\end{enumerate}

\noindent \textbf{Exercise 2.10:}

\begin{enumerate}
  \item [(a)]

        1. To show that H is constant we can just show that it's derivative w.r.t t is zero:
        \begin{align*}
          \frac{d}{dt} H(x'(t), x(t)) = 0
        \end{align*}

        2. Let y(t) = x'(t), thus with chain rule we have
        \begin{align*}
          \frac{d}{dt} H(x'(t), x(t)) & = \frac{d}{dt} [\frac{1}{2}(x'(t))^2 + V(x(t))] \\
                                      & = x'(t)x''(t) + V'(x(t))x'(t)                   \\
                                      & = x'(t) [x''(t) + V'(x(t))]
        \end{align*}

        3. After factoring out the x'(t), we can just substitute the definitions from (2.47)
        \begin{align*}
          \frac{d}{dt} H(x'(t), x(t)) & = x'(t) [F(x(t)) + V'(x(t))] \\
                                      & = x'(t) [-V'x(t) + V'(x(t))] \\
                                      & = x'(t) [0]                  \\
                                      & = 0
        \end{align*}

        4. Since $H'(x'(t), x(t)) = 0$, we therefore conclude that $H(x'(t),x(t))$ is constant for all t.

  \item [(b)]

        1. Substituting these initial conditions into H, we have
        \begin{align*}
          H(v_0, x_0) & = \frac{1}{2}(v_0)^2+x_0 \\
                      & = \frac{1}{2}(1) + 0     \\
                      & = \frac{1}{2}
        \end{align*}

        2. Then we substitute this back into the x'(t) we found in the last part
        \begin{align*}
          x' & = \pm{} \sqrt{H(v_0, x_0) - V(x)}           \\
             & = \pm{} \sqrt{\frac{1}{2} - \frac{1}{2}x^2} \\
             & = \pm{} \sqrt{\frac{1}{2}(1-x^2)}           \\
        \end{align*}

        3. We can observe that the function has equilibrium points at
        \[
          x = \pm{1}
        \]

        4. Thus we can describe the solutions as first starting at x=1, then resting for $t>=0$ time, then moving to -1 and resting for $t>=0$ time, and so on, back and forth. The duration of the rests can infinitely chosen, therefore the number of solutions is also infinite.

        5. To check if the function is twice continuously differentiable, we use (2.47)
        \begin{align*}
          x''(t) & = \frac{d}{dx} V(x)            \\
                 & = \frac{d}{dx}(\frac{1}{2}x^2) \\
                 & = -x
        \end{align*}

        6. When transitioning from resting to moving x''(t) would have to jump from 0 to -1 which is non-continuous. Therefore we conclude there is only a single solution which is twice continuously differentiable and that is when there is no resting at the the bounds of the interval.

  \item [(c)]

        1. First, we can solve for x'(t) at the initial conditions we're given
        \begin{align*}
          (x'(0))^2 + 1^4 & = 1 \\
          x'(0)           & = 0
        \end{align*}

        2. Now, we arrange the equation to solve for x'(t)
        \begin{align*}
          (x'(t))^2 + x^4(t) & = 1            \\
          (x'(t))^2          & = 1-x^4        \\
          x'(t)              & = \sqrt{1-x^4}
        \end{align*}

        3. Like the previous equation, this equation also has equilibrium points at -1 and 1. It behaves the same as the previous equation as well, where it rests at x=1 and x=-1 before moving, thus it also has infinitely many solutions.

        4. To determine if the solutions are twice continuously differentiable, we can differentiate the original equation
        \begin{align*}
          \frac{d}{dt} [(x'(t))^2 + x^4(t)] & = 0 \\
          2x'(t)x''(t) + 4x'(t)x^3(t)       & = 0 \\
          2x'(t) [x''(t) + 2x^3(t)]         & =0
        \end{align*}

        5. Here, we find that when $x'(t) \neq{0}$, we have that
        \[
          x''(t) = -2x^3(t)
        \]

        6. Similar to the last equation, this is the only solution that is twice continuously differentiable because it doesn't rest at either boundaries: -1 and 1.
\end{enumerate}


\noindent \textbf{Exercise 3.2:}
\begin{enumerate}
  \item [(a)]

        1. First we differentiate v(x,y) with chain rule
        \begin{align*}
          v'(x,y) = \frac{x'x-y'y}{v}
        \end{align*}

        2. Then we substitute the x' and y' from $\mathbf{v}(x,y)$ and $v^2 = x^2-y^2$
        \begin{align*}
          v'(x,y) & = \frac{(x-y\sqrt{x^2-y^2})x-(y-x\sqrt{x^2-y^2})y}{v} \\
                  & = \frac{x^2-xyv-y^2+xyv}{v}                           \\
                  & = \frac{x^2-y^2}{v}                                   \\
                  & = \frac{v^2}{v}                                       \\
                  & = v
        \end{align*}

        3. Thus we have that
        \[
          v'(x,y) = v(x,y)
        \]

        4. Next, we differentiate u(x,y) with quotient rule
        \begin{align*}
          u'(x, y) = \frac{1}{2}(\frac{(x-y)(x'+y')-(x+y)(x'-y')}{(x-y)^2})
        \end{align*}

        5. We can simplify this by solving for
        \begin{align*}
          (x' + y') & = (1-v)(x+y) \\
          (x' - y') & = (1-v)(x-y)
        \end{align*}

        6. Substituting back into u'(x,y) we get
        \begin{align*}
          u'(x, y) & = \frac{1}{2}(\frac{(x-y)(x+y)(1-v)-(x+y)(x-y)(1+v)}{(x-y)^2}) \\
                   & = \frac{1}{2}(\frac{(x+y)(1-v-1-v)}{x-y})                      \\
                   & = -v(\frac{x+y}{x-y})
        \end{align*}

        7. And from $2u(x,y) = \frac{x+y}{x-y}$, we have
        \begin{align*}
          u'(x, y) = -2uv
        \end{align*}

        8. There we conclude that our final vector field is
        \begin{align*}
          \mathbf{w}(\mathbf{u}(t)) = \begin{bmatrix} -2uv \\ v \end{bmatrix}
        \end{align*}

  \item [(b)]

        1. We can solve for v(t) first, which is a standard differential equation
        \begin{align*}
          v(t) = v_0e^t
        \end{align*}

        2. Now we can substitute this in to solve for u(t) with separation of variables
        \begin{align*}
          u'(t)              & = -2uv_0e^t         \\
          \int{}\frac{du}{u} & = \int{}-2v_0e^t dt \\
          \ln(u)             & = -2v_0e^t + C      \\
          u(t)               & = Ce^{-2v_0e^t}
        \end{align*}

        3. Now we plug in initial conditions $u(0) = u_0$ and solve for the constant
        \begin{align*}
          u_0 & = Ce^{-2v_0}  \\
          C   & = u_0e^{2v_0}
        \end{align*}

        4. Thus, we have the particular solution for u(t) as well
        \begin{align*}
          u(t) & = u_0e^{2v_0}e^{-2v_0e^t} = u_0e^{2v_0(1-e^t)}
        \end{align*}

        5. Therefore, the general solution is
        \begin{align*}
          \mathbf{u}(t) = \begin{bmatrix}   u_0e^{2v_0(1-e^t)} \\ v_0e^t       \end{bmatrix}
        \end{align*}

  \item[(c)]

        1. We have the invert the coordinates back to x,y. We can simplify this by using the following equations
        \begin{align*}
          x(t) & = \frac{(x+y)+(x-y)}{2} \\
          y(t) & = \frac{(x+y)-(x-y)}{2}
        \end{align*}

        2. Now, we only need to solve for $s = (x+y)$ and $d= (x-y)$
        \begin{itemize}
          \item From definition of u(t):
                \begin{align*}
                  u(t) & = \frac{s}{2d} \\
                  2ud  & = s
                \end{align*}
          \item  From definition of v(t):
                \begin{align*}
                  v(t) & = \sqrt{sd} \\
                  v^2  & = sd
                \end{align*}

          \item  Solve for d:
                \begin{align*}
                  v^2            & = sd                 \\
                  v^2            & = 2ud^2              \\
                  \frac{v^2}{2u} & = d^2                \\
                  d              & =\frac{v}{\sqrt{2u}}
                \end{align*}

          \item Solve for s:
                \begin{align*}
                  s & = 2ud                   \\
                    & = 2\frac{uv}{\sqrt{2u}} \\
                    & = v\sqrt{2u}
                \end{align*}
        \end{itemize}

        3. Then, we plug in the solution we got in the previous part
        \begin{itemize}
          \item For s = x(t) + y(t)
                \begin{align*}
                  x(t) + y(t) & = v(\sqrt{2u})                               \\
                              & = v_0e^t\sqrt{2u_0e^{2v_0(1-e^t)}}           \\
                              & = (x_0 + y_0)e^{t+\sqrt{x_0^2-y_0^2}(1-e^t)}
                \end{align*}
          \item  For d = x(t) - y(t)
                \begin{align*}
                  x(t) - y(t) & = \frac{v}{\sqrt{2u}}                        \\
                              & = \frac{v_0e^t}{\sqrt{2u_0e^{2v_0(1-e^t)}}}  \\
                              & = (x_0 - y_0)e^{t-\sqrt{x_0^2-y_0^2}(1-e^t)}
                \end{align*}
        \end{itemize}

        4. Therefore, we have that the general solution is equal to
        \begin{align*}
          \mathbf{x}(t) = \frac{1}{2} \begin{bmatrix}   (x_0 + y_0)e^{t+\sqrt{x_0^2-y_0^2}(1-e^t)} + (x_0 - y_0)e^{t-\sqrt{x_0^2-y_0^2}(1-e^t)} \\ (x_0 + y_0)e^{t+\sqrt{x_0^2-y_0^2}(1-e^t)} - (x_0 - y_0)e^{t-\sqrt{x_0^2-y_0^2}(1-e^t)}\end{bmatrix}
        \end{align*}

        5. And the flow transformation is
        \begin{align*}
          \Phi_t(\mathbf{x}) =  \frac{1}{2} \begin{bmatrix}   (x + y)e^{t+\sqrt{x^2-y^2}(1-e^t)} + (x - y)e^{t-\sqrt{x^2-y^2}(1-e^t)} \\ (x + y)e^{t+\sqrt{x^2-y^2}(1-e^t)} - (x - y)e^{t-\sqrt{x^2-y^2}(1-e^t)}\end{bmatrix}
        \end{align*}

\end{enumerate}

\noindent \textbf{Exercise 3.4:}

\begin{enumerate}
  \item [(a)]
        1. First we need to find the eigenvalues
        \begin{align*}
          det(A - \lambda{}) & = (-4-\lambda{})(-1-\lambda{}) -10 = 0 \\
                             & = \lambda^{2} -5\lambda{} -6           \\
                             & = (\lambda{}+ 6)(\lambda{} -1)         \\\\
                             & \lambda_{1} = 1,~ \lambda_{2} = -6
        \end{align*}

        2. Now we solve for the eigenvectors
        \begin{enumerate}
          \item For $\lambda_{1}$:
                \begin{align*}
                  \begin{pmatrix} -5 & 2 \\ 5 & -2\end{pmatrix}\mathbf{v}_1 & = \mathbf{0}                           \\
                  \mathbf{v}_1                                              & =  \begin{pmatrix} 2 \\ 5\end{pmatrix}
                \end{align*}
          \item For $\lambda_{2}$:
                \begin{align*}
                  \begin{pmatrix} 2 & 2 \\ 5 & 5\end{pmatrix}\mathbf{v}_2 & = \mathbf{0}                            \\
                  \mathbf{v}_2                                            & =  \begin{pmatrix} 1 \\ -1\end{pmatrix}
                \end{align*}
        \end{enumerate}

        3. Collecting our vectors in matrices we have
        \begin{align*}
          D      & = \begin{bmatrix} 1 & 0 \\ 0 & -6\end{bmatrix}              \\
          V      & = \begin{bmatrix} 2 & -1 \\ 5 &  1\end{bmatrix}             \\
          V^{-1} & = \frac{1}{7} \begin{bmatrix} 1 & 1 \\ -5 &  2\end{bmatrix}
        \end{align*}

        4. Thus, we have that
        \begin{align*}
          x^{tA} & = Ve^{tD}V^{-1}                                                                                                                                                 \\
                 & = \frac{1}{7} \begin{bmatrix} 2 & -1 \\ 5 &  1\end{bmatrix} \begin{bmatrix} e^t & 0 \\ 0 & e^{-6t} \end{bmatrix}  \begin{bmatrix} 1 & 1 \\ 5 &  -2\end{bmatrix} \\
                 & = \frac{1}{7} \begin{bmatrix} 2e^t + 5e^{-6t} & 2e^t - 2e^{-6t} \\ 5e^t-5e^{-6t} &  5e^t+2e^{-6t} \end{bmatrix}
        \end{align*}

        5. And the general solution
        \[
          \mathbf{x}(t) = \frac{1}{7} \begin{bmatrix} 2e^t + 5e^{-6t} & 2e^t - 2e^{-6t} \\ 5e^t-5e^{-6t} &  5e^t+2e^{-6t} \end{bmatrix} \mathbf{x_0}
        \]

  \item [(b)]

        1. The general solution can also be written as the linear combination of the eigenvectors
        \[
          \mathbf{x}(t) = c_1e^{\lambda_{1}t}\mathbf{v}_1 + c_2e^{\lambda_{2}t}\mathbf{v}_2
        \]

        2. Notice that as t goes to infinity, the first term will blow up while the second term goes to zero, thus we conclude
        \[
          \mathbf{x}_0 = c_1\begin{pmatrix} 1 \\ -1 \end{pmatrix} | c_1 \in{\mathbb{R}}
        \]

\end{enumerate}

\noindent \textbf{Exercise 3.5:}

\begin{enumerate}
  \item [(a)]
        1. First we need to find the eigenvalues
        \begin{align*}
          det(A - \lambda{}) & = (5-\lambda{})(1-\lambda{}) -4 = 0 \\
                             & = \lambda^{2} -6\lambda{} + 9       \\
                             & = (\lambda{}_3)^2                   \\\\
                             & \lambda_{1} = \lambda_{2} = 3
        \end{align*}

        2. This matrix is not diagonalizable so we must use the formula for generalized eigenvectors, first we check the conditions
        \begin{align*}
          A - 3I     & = \begin{pmatrix} 2 & -1 \\ 4 & -2 \end{pmatrix}                                                    \\
          (A - 3I)^2 & =  \begin{pmatrix} 2 & -1 \\ 4 & -2 \end{pmatrix}^2 =  \begin{pmatrix} 0 & 0 \\ 0 & 0 \end{pmatrix}
        \end{align*}

        3. Thus every non-zero vector in $\mathbb{R}^2$ is a generalized eigenvectors for eigenvalue 3 and we have the general solution
        \begin{align*}
          e^{tA} & = e^{3t}(\sum_{k=0}^1\frac{t^k}{k!}(A-\lambda{}I)^k)           \\
                 & = e^{3t}(I + t\begin{pmatrix} 2 & -1 \\ 4 & -2 \end{pmatrix})  \\
                 & = e^{3t}\begin{pmatrix} 1 + 2t & -t \\ 4t & 1-2t \end{pmatrix}
        \end{align*}

  \item [(b)]

        1. Because the eigenvalue is three the solution always blows up, therefore the only solution is
        \[
          \mathbf{x}_0 = \mathbf{0}
        \]

\end{enumerate}

\noindent \textbf{Exercise 3.6:}

\begin{enumerate}
  \item [(a)]

        1. First, we solve for the eigenvalues
        \begin{align*}
          det(A - \lambda{}) & = (-4-\lambda{})(-1-\lambda{}) -18 = 0 \\
                             & = \lambda^{2} + 5\lambda{} -14         \\
                             & = (\lambda+7)(\lambda-2)               \\\\
                             & \lambda_{1} = -7,~ \lambda_{2} = 2
        \end{align*}

        2. Now we solve for the eigenvectors
        \begin{itemize}
          \item For $\lambda_{1}$:
                \begin{align*}
                  \begin{bmatrix} 3 & 2 \\ 9 & 6 \end{bmatrix} \mathbf{v}_1 & = \mathbf{0}                            \\
                  \mathbf{v}_1                                              & = \begin{bmatrix} -2 \\ 3 \end{bmatrix}
                \end{align*}
          \item  For $\lambda_{2}$:
                \begin{align*}
                  \begin{bmatrix} -6 & 2 \\ 9 & -3 \end{bmatrix}\mathbf{v}_2 & = \mathbf{0}                           \\
                  \mathbf{v}_2                                               & = \begin{bmatrix} 1 \\ 3 \end{bmatrix}
                \end{align*}
        \end{itemize}

        3. Diagonalizing A we have
        \begin{align*}
          A = VDV^{-1} = \begin{bmatrix} -2 & 1 \\ 3 & 3\end{bmatrix} \begin{bmatrix} -7 & 0 \\ 0 & 2 \end{bmatrix}\frac{1}{9}\begin{bmatrix} -3 & 1 \\ 3 & 2 \end{bmatrix}
        \end{align*}

        4. Thus the general solution is
        \begin{align*}
          e^{tA} & = Ve^{tD}V^{-1}                                                                                                                                                 \\
                 & = \frac{1}{9} \begin{bmatrix} -2 & 1 \\ 3 & 3\end{bmatrix} \begin{bmatrix} e^{-7t} & 0 \\ 0 & e^{2t} \end{bmatrix}\begin{bmatrix} -3 & 1 \\ 3 & 2 \end{bmatrix} \\
                 & = \frac{1}{9} \begin{bmatrix} 3e^{2t} +6e^{-7t} & 2e^{2t}-2e^{-7t} \\ 9e^{2t}-9e^{-7t} & 6e^{2t}+3e^{-7t} \end{bmatrix}
        \end{align*}

  \item [(b)]

        1. The general solution can be rewritten as
        \[
          \mathbf{x}(t) = c_1e^{-7t}\mathbf{v}_1 + c_2e^{2t}\mathbf{v}_2 ~|~ c_1, c_2 \in{\mathbb{R}}
        \]

        2. The second term blows up as t approaches infinity, thus we have that
        \[
          \mathbf{x}_0 = c_1\begin{bmatrix} -2 \\ 3\end{bmatrix} ~|~ c_1 \in{\mathbb{R}}
        \]
\end{enumerate}

\end{document}